\scanpage{28}
\deleted{élémentaire d'Artin \inserted{anabélienne} -- donc son $\pi_1$ \deleted{est} a un centre trivial. Ceci posé.}

\uline{Corollaire 2} La conclusion \inserted{$f=g$} du corollaire précédent \deleted{b)} \deleted{\unsure{sera}} \inserted{est} valable si, \deleted{au lieu de supposer $f$ ou $g$ dominants,} on suppose que \inserted{$k$ est de car. $0$, $X$ \uline{normal}, et} \deleted{enfin:} on est dans l'une des hyp. suivantes

c) l'image de $\pi_1(f)$ est un sous-groupe ouvert de $\pi_1(Y)$, \inserted{$Y$ est une variété élémentaire d'Artin anabélienne}

d) l'image \inserted{de $\pi_1(X)$} par $\pi_1(f)$ \deleted{de tout sous-groupe ouvert de $\pi_1(X)$} a un centralisateur dans $\pi_1(Y)$ réduit à $(1)$, \deleted{\unsure{$X$ normal}} \inserted{et $Y$ se plonge dans un groupe alg. extension d'une VA par un tore}

\deleted{e) Le centre de $\pi_1(X)$, et le centralisateur de son image dans $\pi_1(Y)$, sont réduits à $1$.}
\margin{\fbox{\parbox{0.26\textwidth}{\unsure{Remarquons} si $Y$ est \unsure{simple}, \unsure{centre} $\unsure{\ne1}$ \unsure{réduit} \illegible{3}.}}}

Comme le centralisateur d'un sous-groupe ouvert de $\pi_1(Y)$ ($\pi_1(Y)$ étant ext. successive de groupes profinis \uline{libres} anabéliens) est réduit à $(1)$, comme on a vu${}^*$ plus haut, le cas c) est un cas particulier de d); \deleted{\unsure{de même $X$ aurait}} \deleted{D'autre part, d) se ramène à e), en remplaçant $X$ par un ouvert d'Artin} -- mais dans le cas d), quitte à remplacer $X$ par un ouvert d'Artin (la normalité assurant que $\pi_1(U)\longrightarrow\pi_1(X)$ surj.),

\margin{${}^*$ il faut quand $\unsure{\overline m|=}$ travailler! pour le cas du $\pi_1$ d'une courbe complète de genre $g\geq2$ -- et vérifier que son $\pi_1$ est à centre trivial et qu'il \unsure{n'a pas} d'éléments d'ordre fini! Sinon, se borner dans le cor. au cas de fibres successives \unsure{non complètes}.}
\margin{27}
